\section{Conclusion}
\label{sec:conclusion}

\paragraph{5.1}
The standalone output-capital relation cannot sustain a stable capacity benchmark; cointegration emerges only when the rate of exploitation enters the system. This critical replication of Shaikh (2016) accounting-based methodology demonstrates that the bivariate VECM fails across all lag profiles and deterministic configurations, without any evidence of cointegration between output and capital stocks. Instead, cointegration holds only in a trivariate VECM conditioned on income distribution and historical shock controls. Once trend-containing specifications are excluded as economically inadmissible, the retained models yield a transformation elasticity below unity ($\hat{\theta} < 1.0$), placing the post-war US corporate sector in a structural overaccumulation regime where capital accumulation outpaces capacity formation in the long-run.

\paragraph{5.2}
The non-unitary transformation elasticity imposes a secular downward drag on the rate of profit that capitalists can only counteract by intensifying exploitation or idling capacity. Within the classical-Marxian decomposition $r = (\Pi/Y) \cdot \mu \cdot (Y^p/K)$, a transformation elasticity below unity forces the capital-capacity ratio ($Y^p/K$) into a persistent decline as capital accumulates. The speed of adjustment of the system suggests that capital accumulation is weakly exogenous ($\alpha_k \approx 0$), while the rate of exploitation serves as the primary error-correction channel restoring long-run reproduction. This weak exogeneity of investment supports the Sraffian position over the Neo-Kaleckian demand-endogenous view. Yet the structural overaccumulation ($\hat{\theta} < 1.0$) introduces a counteracting dynamic: the secular instability of unbalanced growth and the weak exogeneity of investment may cancel each other out, obscuring the resolution of this debate. This theoretical impasse reflects the shared balanced-growth assumption underlying both frameworks, which forces a fixed output-capital ratio onto a historically shifting accumulation process. Ultimately, the capacity ceiling is not a neutral engineering constraint but a contingent outcome of decentralized accumulation, shaped by the historical balance of class power and the distributive conflicts of the capitalist labor process.

\paragraph{5.3}
The single-sector bounding and the time-invariant parameterization of $\theta$ represent the two primary analytical limitations of this chapter. By abstracting from inter-sectoral disproportionalities, capacity transfers, and uneven development across departments of production, the single-sector framework isolates the aggregate capital-to-capacity relation at the cost of tracing how sectoral imbalances spill over to affect the general rate of profit. Similarly, treating $\theta$ as a constant over the 1947–2011 period forces the model to average across historically distinct choice-of-technique regimes and social structures of accumulation. While external demand components—military spending, state deficits, and consumer debt expansion—may stabilize the overaccumulation tendency, these transmission mechanisms remain hypotheses rather than findings directly estimated here.

\paragraph{5.4}
The empirical failure of the standalone output-capital relation establishes a clear theoretical imperative: capacity measurement must integrate workplace-level conflict \emph{ex ante} rather than appending it \emph{ex post}. The immediate next step is to endogenize the transformation elasticity $\theta$ as a time-varying parameter driven by class struggle and the choice of technique, tracing how the capital-capacity conversion shifts across historical US accumulation regimes. Building on this closed-economy foundation for the center, the dissertation must subsequently confront the structural realities of the periphery. In peripheral economies, the productive ceiling is not merely bounded by domestic distributive conflict, but is structurally truncated by the Balance of Payments (BoP) constraint. For these economies, potential capacity is subordinated to the external viability of the accumulation process; the imperative to generate sufficient foreign exchange to finance essential imports and service external debt acts as a hard, macroeconomic ceiling on domestic expansion. The next analytical layer will therefore relax the single-sector bounding to incorporate open-economy dynamics, estimating a capacity utilization rate for peripheral countries that explicitly conditions the output-capital relation on these external constraints. This upgrade will reveal whether chronic idle capacity in the periphery is driven by domestic overaccumulation or by the structural necessity of maintaining external equilibrium, ultimately demonstrating that the capacity ceiling in the Global South is a geopolitical and distributive construct rather than a neutral technical limit.

